\documentclass[a4paper]{article}

% Español (México) con XeLaTeX: fontspec para fuentes Unicode, polyglossia
% para la división silábica y los nombres del documento en la variante mexicana.
\usepackage{fontspec}
\usepackage{polyglossia}
\setdefaultlanguage[variant=mexican]{spanish}
% Los nombres chinos van como «pinyin (original)»: xeCJK da a los caracteres
% Han su propia fuente; sin ella XeLaTeX los omite con un aviso.
\usepackage{xeCJK}
\setCJKmainfont{FandolSong-Regular.otf}
% polyglossia no traduce los nombres de listings.
\AtBeginDocument{\providecommand{\lstlistingname}{}\renewcommand{\lstlistingname}{Listado}%
  \providecommand{\lstlistlistingname}{}\renewcommand{\lstlistlistingname}{Índice de listados}}
\usepackage{amsmath, amssymb}
\usepackage{graphicx}
\usepackage[margin=2.5cm]{geometry}
\usepackage[most]{tcolorbox}
\usepackage{etoolbox}
\usepackage{listings}
\usepackage{hyperref}
\usepackage{booktabs}
\usepackage{array}
\usepackage{tabularx}
\usepackage{longtable}
\usepackage{float}
\usepackage{tikz}
\usetikzlibrary{arrows.meta,positioning,calc}

\newcolumntype{Y}{>{\raggedright\arraybackslash}X}

\newtcolorbox{knowledgebox}[1]{
    enhanced, colback=blue!5!white, colframe=blue!75!black, colbacktitle=blue!75!black,
    coltitle=white, fonttitle=\bfseries, title={#1}, attach boxed title to top left={yshift=-2mm, xshift=2mm},
    boxrule=1pt, sharp corners
}

\newtcolorbox{importantbox}[1]{
    enhanced, colback=yellow!10!white, colframe=yellow!80!black, colbacktitle=yellow!80!black,
    coltitle=black, fonttitle=\bfseries, title={#1}, sharp corners
}

\newtcolorbox{warningbox}[1]{
    enhanced, colback=red!5!white, colframe=red!75!black, colbacktitle=red!75!black,
    coltitle=white, fonttitle=\bfseries, title={#1}, sharp corners
}

\lstset{
    language=Python,
    basicstyle=\ttfamily\small,
    keywordstyle=\color{blue},
    stringstyle=\color{red!60!black},
    commentstyle=\color{green!60!black},
    breaklines=true,
    frame=single,
    numbers=left,
    numberstyle=\tiny\color{gray},
    captionpos=b,
    extendedchars=true
}

\newcommand{\notetitle}{Notas de entrevista: la próxima parada de los GUI Agent}
\newcommand{\noteauthors}{Elaborado a partir de Ungrounded 不着边际 EP01}
\newcommand{\notedate}{2026-05-08}
\newcommand{\videochannel}{Ungrounded 不着边际}
\newcommand{\videopublishdate}{2026-01-28}
\newcommand{\videoduration}{02:27:24}
\newcommand{\videourl}{https://www.bilibili.com/video/BV1y36jBnEiD/}
\newcommand{\videocoverpath}{cover.jpg}
\newcommand{\repourl}{https://github.com/hqhq1025/ai-course-notes}

\usepackage{fancyhdr}
\pagestyle{fancy}
\fancyhf{}
\fancyfoot[C]{\thepage}
\fancyfoot[R]{\footnotesize\color{gray}\href{\repourl}{AI Course Notes}}
\renewcommand{\headrulewidth}{0pt}
\renewcommand{\footrulewidth}{0pt}

\begin{document}

\begin{titlepage}
\centering
{\Large Notas de entrevista\par}
\vspace{1.2cm}
{\huge\bfseries \notetitle\par}
\vspace{0.8cm}
{\large \noteauthors\par}
\vspace{0.3cm}
{\large \notedate\par}
\vspace{1.2cm}

\ifdefempty{\videocoverpath}{
{\small Al generar las notas, indique la ruta local de la portada del video.\par}
}{
\includegraphics[width=0.82\textwidth,height=0.45\textheight,keepaspectratio]{\videocoverpath}\par
}

\vfill
\begin{tcolorbox}[width=0.9\textwidth, colback=black!2!white, colframe=black!60, sharp corners]
\textbf{Autor o canal del video}: \videochannel\par
\textbf{Fecha de publicación}: \videopublishdate\par
\textbf{Duración del video}: \videoduration\par
\textbf{Enlace del video}: \href{\videourl}{\nolinkurl{\videourl}}\par
\textbf{Página del proyecto}: \href{\repourl}{\nolinkurl{\repourl}}\par
\end{tcolorbox}
\end{titlepage}

\tableofcontents
\newpage

\section*{Introducción}

El valor de esta conversación no está en ofrecer un «panorama rápido del estado actual de los GUI Agent», sino en desdoblar una y otra vez la misma pregunta desde cuatro niveles: investigación, producto, negocio y organización. Cuando los modelos por fin empiezan a poder operar entornos digitales reales, ¿deberíamos entenderlos como una herramienta de automatización que sabe dar clic en la pantalla, o como un nuevo tipo de agente digital que apenas se está formando? Las trayectorias de Xie Tianbao (谢天宝), Gu Yu (谷雨) y Shang Yanyi (尚晏仪) constituyen justo tres perspectivas complementarias: una parte de OSWorld, GUI grounding y benchmark; otra parte de semantic parsing, Mind2Web y la evaluation académica; y otra parte de la práctica comercial de construir productos de GUI Agent desde 2021 y luego pasar a los viajes con IA. Los tres no ofrecen un lema unificado, pero calibran una y otra vez alrededor del mismo conjunto de tensiones: escenarios reales frente a benchmark, entrada universal frente a producto vertical, operación visible frente a experiencia de usuario, memoria de largo plazo frente a la deriva conceptual, ventaja de velocidad frente a activo acumulable.

Si solo se narrara este episodio en orden cronológico, sería fácil leerlo como una serie de opiniones sueltas del sector: Operator no es fácil de usar, la puntuación en benchmark no importa tanto, los Agent podrían afectar a la publicidad, memory es difícil, un startup necesita un foso competitivo. Pero lo verdaderamente interesante de esta conversación es que estas opiniones no están aisladas. El problema de experiencia de Operator termina remitiendo a «si el usuario en realidad quiere o no ver cómo opera el Agent»; la crisis de benchmark termina remitiendo a «si un producto real puede generar un reward confiable»; la discusión sobre Unified Action Space termina remitiendo a «por qué Deep Research, Coding, Browser Use y la operación de GUI no deberían fragmentarse»; y la advertencia sobre la deriva conceptual termina remitiendo a «cómo mantiene una organización una semántica compartida». Dicho de otro modo, esto no es una lista de funciones sobre los GUI Agent, sino una discusión metodológica sobre «cómo entra un agente digital al mundo real».

Por eso, esta nota no se organiza estrictamente en orden cronológico, sino que se reorganiza en ocho temas. Cada tema procura conservar el contexto concreto de la conversación: desde qué experiencia parte cada quien, por qué llega a ese juicio, qué malentendido común refuta, y qué significa ese juicio para investigadores, equipos de producto o emprendedores, respectivamente.

\begin{importantbox}{Tres hilos conductores para leer esta nota}
El primer hilo es «el problema de la evaluación»: desde OSWorld y Mind2Web hasta los productos reales, qué tipo de evaluation puede realmente guiar el progreso de un Agent. El segundo hilo es «el problema del producto»: desde Operator, el teléfono Doubao y Unified Action Space hasta el context switching, qué tipo de agente digital necesita en realidad el usuario. El tercer hilo es «el problema de la acumulación»: desde memory, knowledge y la biblioteca semántica de una organización hasta el foso competitivo de un startup, cómo forma un sistema de Agent activos de largo plazo en interacciones reales.
\end{importantbox}

Por último, cabe aclarar que el video original no tiene subtítulos oficiales; esta nota se elaboró a partir de una transcripción local con Whisper y una revisión línea por línea. Como el video es principalmente una conversación entre varias personas y no una exposición con diapositivas, el cuerpo del texto no inserta capturas de personas de forma forzada; la portada conserva la portada original del video, y los diagramas se usan sobre todo para condensar relaciones conceptuales, no como decoración.
\section{De «saber operar la computadora» a «entender la tarea digital»}

Las trayectorias de entrada de los tres invitados muestran precisamente que el GUI Agent no es una tendencia nueva que haya surgido de repente. Gu Yu (谷雨) empieza por el semantic parsing: la investigación temprana se centraba en cómo convertir una instruction en lenguaje natural en una consulta a una base de datos, una operación sobre una base de conocimiento u otra formal representation. En esa etapa el entorno era más controlable y la tarea se parecía más a «traducir una frase en una estructura ejecutable». Pero la naturaleza del problema ya se acercaba mucho a la del computer-using agent actual: el usuario da un objetivo en lenguaje natural y el sistema tiene que convertirlo en una acción ejecutable dentro del entorno.

La trayectoria de Xie Tianbao (谢天宝), en cambio, se despliega a lo largo de agent benchmark y entornos GUI reales. Desde AgentBench, intercode y OSWorld hasta el GUI grounding multimodal, lo que le interesa es cómo lograr que el modelo haga cosas de verdad en la pantalla, en páginas web y en el sistema operativo, y no solo responda preguntas en el espacio del texto. Shang Yanyi (尚晏仪) aporta una tercera línea: desde 2021 desarrollando productos capaces de operar la computadora en lugar del usuario, después haciendo Web Agent y GUI Agent de propósito general, hasta descubrir que una entrada de propósito general es demasiado difícil de convertir en producto comercial y virar hacia escenarios concretos como los viajes con IA. Al poner juntas las tres líneas, se ve que lo central del GUI Agent no es el «clic del ratón» en sí, sino cómo se conectan entre sí el lenguaje, la visión, la acción, las preferencias del usuario y los procesos de negocio reales.

\begin{knowledgebox}{Antecedentes históricos del GUI Agent}
\begin{tabularx}{\textwidth}{>{\raggedright\arraybackslash}p{0.26\textwidth} Y}
\toprule
Pista de la historia previa & Su significado para el GUI Agent actual \\
\midrule
Semantic Parsing & Entrenar sistemas para convertir tareas en lenguaje natural en representaciones ejecutables, lo que ofreció una forma temprana de «del lenguaje a la acción». \\
Web Automation / RPA & Demostró que muchas tareas digitales pueden convertirse en flujos de proceso, pero también expuso la fragilidad de los sistemas basados en reglas frente a entornos dinámicos. \\
Agent Benchmark & Usó tareas como OSWorld y Mind2Web para hacer explícito el problema de las capacidades, de modo que los modelos pudieran compararse entre sí. \\
Exploración de conversión en producto & Expuso la distancia entre un demo de investigación y la experiencia del usuario: poder ejecutar no equivale a ser útil, y poder mostrar el proceso no equivale a generar confianza. \\
\bottomrule
\end{tabularx}
\end{knowledgebox}

El cambio de 2025 es que estas líneas finalmente volvieron a conectarse gracias a modelos multimodales más potentes. Xie Tianbao menciona que los modelos tempranos ni siquiera podían realizar bien acciones aparentemente simples como «hacer clic en el botón de cerrar de la esquina superior derecha», «arrastrar el volumen a 80\%» o «modificar el nombre en el campo de entrada». Después, mediante datos sintéticos a gran escala, la expansión de entornos y la iteración de modelos, el GUI grounding pasó de ser «puro concepto» a resultar «básicamente utilizable». Esto no es un ajuste menor, porque la dificultad del entorno GUI no está solo en ver el botón, sino en alinear al mismo tiempo la localización visual, la instrucción en lenguaje, las coordenadas de la pantalla, el tipo de acción y el objetivo de la tarea.

Pero es precisamente aquí donde el problema del producto empieza a agudizarse. La actitud de Shang Yanyi al hablar de Operator es representativa: le hace sentir a quienes trabajan en el área que la dirección ha sido validada, pero ella no considera que sea un buen producto. Ver a una pequeña computadora hacer clic lentamente en una página web confunde con facilidad al usuario; una acción que él mismo podría completar en unos segundos, mostrada por la IA a una velocidad muy lenta, no necesariamente aumenta la confianza y puede, por el contrario, exponer torpeza. El investigador quiere demostrar que «de verdad operé por ti», mientras que el usuario tal vez solo quiere saber si «el resultado es confiable y si necesito confirmarlo».

\begin{warningbox}{Equiparar el GUI Agent con «ver la pantalla y hacer clic» hace perder de vista el problema real}
Las acciones superficiales del GUI Agent son hacer clic, escribir, desplazar y arrastrar; pero su problema profundo es la comprensión de la tarea digital. Un Agent maduro debería saber cuándo hacer clic en la GUI, cuándo llamar a una API, cuándo escribir un script, cuándo leer el archivo directamente y cuándo pedirle confirmación al usuario. La GUI es una parte del espacio de acción, no la inteligencia en sí misma.
\end{warningbox}

Lo que este capítulo establece en realidad es la premisa subyacente de toda la conversación: el futuro del GUI Agent no es solo un grounding más fuerte, ni tampoco una «pequeña computadora remota» más realista. Lo que tiene que resolver es la comprensión del objetivo, la selección de la acción, el aprendizaje a partir de la retroalimentación y los límites de la confianza dentro del entorno digital. Solo si esta premisa queda establecida, las discusiones posteriores sobre benchmark, modelo de negocio, UDA, memory y foso competitivo podrán conectarse en una sola línea.

\subsection{Resumen de la sección}

El GUI Agent no es un concepto nuevo que haya surgido de la nada, sino el resultado de la convergencia de semantic parsing, web automation, agent benchmark y la exploración de conversión en producto. La capacidad multimodal de 2025 hizo que «operar la computadora» entrara en un rango utilizable, pero lo verdaderamente difícil es pasar de «ejecutar la acción» a «entender la tarea digital». La controversia en torno a productos como Operator muestra que todavía hay una distancia considerable entre la demostración de investigación y la experiencia del usuario.
\section{La crisis de los benchmarks: de puntuaciones comparables a señales de escenarios reales}

El segmento más sólido de esta conversación es el que trata sobre benchmark y evaluation. El punto de partida de Gu Yu (谷雨) es claro: lo que la academia y la industria necesitan de la evaluation nunca ha sido lo mismo. La academia quiere experimentos controlables para verificar hipótesis, definir problemas nuevos y comparar la capacidad de los métodos; la industria, en cambio, se preocupa más por si la mejora en la evaluación puede transferirse al despliegue real, si de verdad permite que el usuario complete la tarea, quede satisfecho, se quede o pague.

Esta distinción importa porque, en este momento, muchas personas dicen por un lado que la puntuación de benchmark no importa, y por otro lado no pueden prescindir del benchmark. Xie Tianbao (谢天宝) ofrece un juicio más matizado: el benchmark no es inútil, sino que depende de en qué etapa se encuentre. En la etapa de 0 a 1, puede definir el problema, exponer las carencias y demostrar que el modelo actual no lo logra; una vez que se entra en la etapa de optimización de 60 a 99, la puntuación pública cada vez es más susceptible a la síntesis de datos, al sobreajuste del conjunto de entrenamiento y a la personalización de la tarea. En ese momento, mirar solo la puntuación puede llevar a un juicio erróneo sobre la utilidad real.

\begin{importantbox}{El cambio de lugar del benchmark}
El benchmark funciona mejor como signal, no como final answer. Puede decirle a la comunidad «aquí hay un problema»«el modelo tiene una carencia aquí»«cierta capacidad está empezando a surgir»; pero no puede responder por sí solo «si este Agent es útil en un producto real». Cuanto más cerca se está de la implementación en la industria, más se necesitan usuarios reales, APIs reales, casos reales y un ciclo continuo de retroalimentación para calibrar.
\end{importantbox}

Shang Yanyi (尚晏仪) habla del reward a partir de los productos de viajes y así completa el eslabón más difícil para la industria. En un producto de viajes, al final se puede observar «si se vendió un boleto o no», pero ese reward está demasiado lejos de cada decisión que toma el modelo paso a paso. Que el usuario no compre el boleto puede deberse a que el precio no le convenció, a que no le gustó la ruta, a que no confía en la IA, o simplemente a que cambió de planes en el último momento. Que el usuario diga «está bien» tampoco equivale a que en verdad vaya a comprar. Es decir, el resultado final es muy claro, pero la atribución del proceso es muy difusa; hay mucha retroalimentación subjetiva, pero pocas señales entrenables.

Esto es distinto de la dificultad de los benchmark académicos. En las tareas académicas, el reward suele estar definido de antemano: se suman puntos si se acierta y se restan si se falla. Pero en un producto real, ni siquiera el usuario sabe con certeza qué experiencia es una buena experiencia. Por ejemplo, cuándo debe interrumpir un asistente de viajes por voz, cuándo debe escuchar, cuándo debe dar una explicación extensa: estos problemas no se pueden resolver con que un investigador escriba un benchmark guiado por su intuición. Que Shang Yanyi diga que dependen más de A/B Test y de una infraestructura de experimentación sostenible es, en realidad, reconocer que la preferencia del usuario no se escribe, sino que se revela poco a poco en la entrega real.

\begin{knowledgebox}{0 a 1 y 60 a 99}
\begin{tabularx}{\textwidth}{>{\raggedright\arraybackslash}p{0.2\textwidth} Y Y}
\toprule
Etapa & Valor del benchmark & Riesgo principal \\
\midrule
0 a 1 & Define tareas nuevas, demuestra que el modelo no puede hacerlas y construye un lenguaje común para el campo. & La tarea puede desligarse del escenario real y convertirse en «benchmark por el benchmark». \\
60 a 99 & Funciona como señal de una capacidad local y ayuda al equipo de ingeniería a ubicar sus carencias. & La puntuación se contamina fácilmente con el entrenamiento, el aumento de datos, la personalización de la tarea y la presentación selectiva. \\
Producto real & Debe usarse junto con la retroalimentación del usuario, los resultados comerciales y la atribución de fallas. & El reward es disperso, subjetivo y de cadena larga, y es difícil convertirlo directamente en señal de entrenamiento. \\
\bottomrule
\end{tabularx}
\end{knowledgebox}

El problema más amplio detrás de este segmento es que la investigación en Agent está pasando de «si la capacidad existe» a «cómo la capacidad produce valor de manera estable». En el primer caso, el benchmark es un instrumento científico; en el segundo, el benchmark es solo una parte del sistema de operación del producto. Lo verdaderamente difícil no es obtener una puntuación, sino conectar en un ciclo la puntuación, la retroalimentación del usuario, el análisis humano, el A/B Test, la selección del modelo, la estrategia de prompt, el flujo de herramientas y las métricas del negocio.

\begin{warningbox}{No atribuyas el fracaso de un producto a la puntuación de un solo modelo}
Cuando un producto de Agent fracasa, es muy probable que no sea algo tan simple como que «el modelo no es lo bastante fuerte». Tal vez la tarea se definió mal, tal vez el reward está demasiado lejos, tal vez la interacción no es confiable, o tal vez el usuario simplemente no quiere completar la tarea de esa manera. El benchmark solo puede iluminar una parte de las causas; no puede sustituir la atribución del producto.
\end{warningbox}

\subsection{Resumen de la sección}

La esencia de la crisis del benchmark no es que la puntuación pierda sentido, sino que se estrecha el margen de su interpretación. La academia todavía necesita el benchmark para definir problemas, mientras que la industria debe colocar la evaluation dentro del ciclo del producto real. La evaluación de los GUI Agent es especialmente difícil porque enfrenta tareas de cadena larga, experiencias subjetivas y reward disperso; la infrastructure que de verdad tiene valor es la infrastructure capaz de probar sin cesar, atribuir sin cesar y convertir sin cesar la retroalimentación del usuario en mejoras del sistema.

\section{De la puerta de entrada publicitaria al trabajo intelectual: una revaloración del modelo de negocio del Agent}

La discusión sobre el modelo de negocio empieza con una comparación muy clásica: los grandes modelos de negocio de la era de internet son sobre todo publicidad, suscripción y comisión por transacción. La publicidad vende atención, la suscripción vende eficiencia de herramienta, la comisión vende intermediación de transacciones. Cuando el Agent funciona como puerta de entrada, desafía directamente a la publicidad, porque el sistema publicitario quiere que el usuario permanezca, navegue y haga clic, mientras que un buen Agent debería precisamente reducir el tiempo que el usuario permanece, y hacer por él el filtrado, la comparación y la ejecución.

El juicio de Shang Yanyi (尚晏仪) resulta muy revelador: lo que completa internet es la conexión de información, y el gasto publicitario corresponde a la distribución de información; lo que produce el Agent es trabajo intelectual, y en la sociedad humana el trabajo intelectual suele pagarse mediante salario, cuota de servicio, comisión de intermediación u honorario. Por eso ella confía más en la suscripción y la comisión que en forzar al Agent de vuelta a la lógica publicitaria. Si el Agent en verdad funciona como secretario, asesor o representante, su valor comercial debería provenir de ahorrar tiempo, reducir el costo de decisión, intermediar transacciones y asumir el resultado de la acción, y no de hacer que el usuario vea el contenido unos segundos más.

El escenario de los viajes funciona en su práctica precisamente porque en él existe de manera natural el papel histórico del travel agent. El usuario entiende qué significa un «representante» y también acepta pagar por planeación, filtrado, transacción y servicio. Los viajes además tienen un ciclo de transacción claro: hotel, vuelo, ruta, visa y servicios de itinerario pueden conectarse a una comisión o a una cuota de servicio. Esto es distinto de una «superpuerta de entrada» generalizada, que parece tener más espacio pero cuyo reward, cumplimiento y ciclo comercial son más difusos.

\begin{importantbox}{El Agent se parece más a una fuerza de trabajo que a un espacio publicitario}
Si el valor de un sistema consiste en pensar, filtrar, actuar y asumir el costo del proceso en lugar del usuario, entonces se acerca más al trabajo intelectual. La lógica publicitaria exige ocupar la atención del usuario; la lógica del Agent exige liberar la atención del usuario. Las dos no son incompatibles, pero el centro de negocio por defecto es distinto.
\end{importantbox}

La discusión sobre el teléfono Doubao lleva este asunto a un nivel más alto. La ventaja del Agent a nivel de teléfono es muy clara: el teléfono conoce la dirección, las preferencias, las cuentas, las conversaciones, la agenda y los hábitos de vida del usuario, por lo que tiene la mejor oportunidad de convertirse en secretario de vida. Pero al mismo tiempo toca la estructura de intereses existente en la distribución de aplicaciones, la recomendación publicitaria y las puertas de entrada a las transacciones. Dicho de otro modo, el Agent de teléfono es técnicamente muy razonable, pero comercialmente cambiará el sustento de muchas personas.

Aquí hay un detalle que vale la pena retener. Shang Yanyi dice que si el usuario le dice directamente al Agent «quiero ir a una isla», esto es más contundente que un sistema de publicidad de video corto que observa que el usuario se detiene unos segundos más en videos de islas. Porque lo primero es una intención explícita, mientras que lo segundo es solo una inferencia de conducta. Pero la intención explícita no equivale automáticamente a éxito comercial. El Agent todavía tiene que estructurar la intención, encontrar la oferta, comparar opciones, construir confianza y completar la transacción. El sistema publicitario es hábil para convertir conductas ambiguas en recomendaciones; el sistema de Agent, en cambio, tiene que convertir lenguaje explícito en acción confiable.

\begin{knowledgebox}{De la economía de la atención a la economía del representante}
\begin{tabularx}{\textwidth}{>{\raggedright\arraybackslash}p{0.2\textwidth} Y Y}
\toprule
Modelo & Origen del valor & Relación con el Agent \\
\midrule
Publicidad & Permanencia, navegación, clic y exposición del usuario. & Un buen Agent reduce la permanencia, por lo que hay un conflicto natural. \\
Suscripción & El usuario paga por eficiencia, capacidad de herramienta y confiabilidad. & Apto para escenarios de alta frecuencia, explícitos y de mejora de eficiencia sostenible. \\
Comisión & La plataforma intermedia la transacción y participa en el cumplimiento. & Apto para viajes, compras, comercio electrónico, reclutamiento y otros escenarios con ciclo de transacción. \\
Servicio de representación & El usuario paga por trabajo intelectual, filtrado, ejecución y límites de responsabilidad. & La forma de largo plazo más cercana a un Agent fuerte. \\
\bottomrule
\end{tabularx}
\end{knowledgebox}

\subsection{Resumen de la sección}

El producto Agent no puede simplemente replicar el modelo publicitario de internet. Su valor se parece más al trabajo intelectual, al servicio de representación y a la intermediación de transacciones. Los viajes son adecuados como caso precisamente porque tienen tanto un papel tradicional de representante como un ciclo de transacción claro. El potencial del Agent de teléfono proviene del contexto personal, pero lo verdaderamente difícil es convertir ese contexto en acción confiable, mientras se manejan los conflictos comerciales con el ecosistema existente.
\section{Forma del producto: de Operator a UDA}

Si el modelo de negocio responde «cómo gana dinero el Agent», la forma del producto responde «qué es lo que el usuario realmente quiere usar». En este punto, los tres invitados no imaginaron al GUI Agent como un robot que permanece siempre visible en la pantalla. Gu Yu quiere una secretaria de vida, capaz de encargarse de visas, agenda, viajes y asuntos varios; Shang Yanyi (尚晏仪) quiere algo más parecido a un coach de vida o a un espejo, que combine contexto de largo plazo para ayudarla a organizar información, hacer las cosas y darle retroalimentación. El punto en común es que a ninguno le interesa ver al Agent hacer clic paso a paso; ambos quieren que el Agent organice sus acciones alrededor de un objetivo.

Este es precisamente el contexto en el que aparecen UDA y el Unified Action Space. Que Deep Research, Coding, Browser Use y GUI Agent se conviertan en productos distintos es resultado de la división del trabajo actual en la industria, pero el objetivo real del usuario no se divide según esas fronteras de producto. Investigar un problema puede requerir buscar en páginas web, leer PDF, escribir código, organizar archivos y enviar mensajes; reservar un viaje puede requerir leer el historial de chats, consultar la visa, comparar hoteles, usar mapas y abrir una agencia de viajes en línea. El objetivo del usuario ocurre dentro del mismo entorno digital, pero el producto exige que el usuario traslade el contexto entre varios Agentes; esto claramente no es el estado final.

\begin{importantbox}{El verdadero significado del Unified Action Space}
El Unified Action Space no significa «poner todos los botones dentro de una superbarra de herramientas», sino permitir que el Agent elija la manera de actuar según el objetivo. Si se puede llamar a una API, que no simule clics; si se puede escribir un script, que no opere página por página; si se puede leer un archivo directamente, que no use captura de pantalla con OCR; solo se entra a la GUI cuando es necesario. Lo que se unifica no es la interfaz, sino la capacidad de organizar acciones alrededor del objetivo del usuario.
\end{importantbox}

Esto también explica por qué «Computer-Using Agent» es más preciso que «GUI Agent». El uso de la GUI es un subconjunto del uso de la computadora. Atar la capacidad del Agent al clic visual hace que el sistema se vuelva deliberadamente torpe en muchos escenarios. Un agente digital verdaderamente maduro debe poder cambiar entre API, código, archivos, navegador, GUI, mensajes y confirmación del usuario, en lugar de traducir todas las tareas a operaciones de mouse.

La discusión sobre la interfaz generativa completa la otra mitad. Shang Yanyi no ve con buenos ojos que «toda la interfaz se genere de manera instantánea por el modelo», porque el usuario necesita estabilidad y previsibilidad. La interfaz no es solo un contenedor de información; también es una estructura de hábito y confianza del usuario. Si cada vez cambian la posición de los botones, la forma de interactuar y la presentación de los resultados, el usuario se sentirá inseguro y le costará adquirir soltura. Una dirección más realista quizá sea combinar componentes estables con una organización dinámica: las operaciones básicas, el proceso de confirmación y los límites de permisos se mantienen estables, mientras que la presentación concreta de los resultados y la asistencia de la tarea pueden generarse de manera dinámica.

\begin{warningbox}{El error común de la interfaz generativa}
El valor de la interfaz generativa no es hacer que la interfaz resulte extraña cada vez, sino acercarla más a la tarea actual. Si se sacrifica la previsibilidad para perseguir una «sensación de IA», el producto termina siendo más difícil de usar. Para el Agent, los componentes semánticos estables importan más que una generación instantánea vistosa.
\end{warningbox}

Aquí hay además un problema que la versión anterior del resumen subestimó: la privacidad y la confianza. Mientras más útil sea un Agent tipo secretaria de vida, más necesitará acceder a domicilio, cuentas, chats, agenda, preferencias, pagos e información de identidad. Puede ser local, puede usar aprendizaje federado, puede establecer un límite de confianza como el de un administrador de contraseñas, pero, sea cual sea la ruta, debe hacer saber al usuario qué se recuerda, qué se puede eliminar y qué acción requiere confirmación. La forma del producto del Agent no es solo una cuestión de forma de la interfaz; en esencia, es una forma de permisos.

\subsection{Resumen de la sección}

De Operator a UDA, lo que cambia es el centro del producto: de «mostrar cómo opera la IA» a «organizar acciones alrededor del objetivo del usuario». El Unified Action Space busca resolver la fragmentación de capacidades y el traslado de contexto; el Computer-Using Agent busca evitar que todas las tareas se degraden a clics de GUI; la interfaz generativa debe servir a la tarea, no romper la previsibilidad. Un producto de Agent verdaderamente maduro debe manejar al mismo tiempo la elección de acciones, la confirmación del usuario y los límites de privacidad.
\section{La paradoja de las tareas largas: entre más pueda correr el Agent, más se cansa el usuario}

La observación más cercana a la experiencia cotidiana en este episodio es la del context switching. Shang Yanyi (尚晏仪) dice que, después de usar mucho Deep Research y coding agents, siente una carga muy fuerte: se lanza un prompt completo y el Agent corre durante mucho tiempo; si la capacidad es mala, ella no lo usa; si la capacidad es buena, tampoco necesita ver la salida intermedia. Entonces no sabe qué hacer mientras espera, así que cambia a otra tarea y regresa unos minutos después a revisar, lo cual la cansa aún más.

Esta observación es muy importante porque refuta buena parte de la narrativa por defecto de muchos productos de Agent. Hoy varios productos usan «poder hacer una tarea de una hora» como argumento de venta, pero Gu Yu (谷雨) señala que si se pudiera lograr la misma calidad en 30 segundos, nadie querría esperar una hora. La capacidad de sostener tareas largas sin duda tiene valor, pero es una capacidad del sistema, no la experiencia del usuario en sí misma. Presentar la lentitud como autonomía es una narrativa de producto muy peligrosa.

\begin{importantbox}{La capacidad de tareas largas no es una experiencia de espera larga}
Lo que el usuario quiere es que la tarea se complete de forma confiable, no verse obligado a administrar una caja negra que corre durante mucho tiempo. Que el Agent pueda correr mucho tiempo indica que quizás posee capacidad de ejecución autónoma; pero el valor del producto viene de menos espera, menos revisión, menos cambio de contexto cognitivo y una mayor calidad de resultado.
\end{importantbox}

La conversación también divide este problema en capa de producto y capa de capacidad. La capa de producto puede mejorar la sensación mediante estados intermedios, retroalimentación de progreso, mecanismos de interrupción y puntos de confirmación clave. Por ejemplo, algunos coding agents le informan continuamente al usuario qué están haciendo, mientras que otros sistemas pasan varios minutos sin dar salida alguna y eso genera ansiedad en el usuario. Pero la capa de capacidad es más fundamental: que el Agent busque información con lentitud, use herramientas sin soltura, siga rutas de decisión enredadas y abra repetidamente páginas web inútiles indica que todavía no posee esa capacidad de juzgar rápidamente el valor de la información que tiene un humano experimentado.

Esto también explica por qué «escribir un prompt para que el Agent corra seis horas» solo le sirve a una minoría de usuarios. Las tareas reales suelen necesitar aclaraciones e iteraciones de ida y vuelta. Ni siquiera el propio usuario puede necesariamente escribir de una vez un requerimiento completo. Si un Agent toma un requerimiento ambiguo y corre durante mucho tiempo, es muy probable que al final produzca una gran cantidad de trabajo que parece esforzado pero que en realidad se desvió. La autonomía real no debería ser «correr solo durante mucho tiempo», sino «saber cuándo avanzar, cuándo detenerse y cuándo pedirle un juicio al usuario».

\begin{knowledgebox}{Cuatro puntos de diseño para la experiencia de tareas largas}
\begin{itemize}
  \item Compresión de estado: decirle al usuario qué problema se está resolviendo en este momento, en lugar de volcarle los registros de bajo nivel.
  \item Puntos de verificación: pedir confirmación en las bifurcaciones de alto riesgo, en lugar de entregar una sorpresa de una sola vez al final.
  \item Interrumpibilidad: permitir que el usuario cambie de rumbo, pause o agregue información.
  \item Resumen de diferencias: al terminar, explicar qué se hizo, qué se cambió y qué sigue siendo incierto.
\end{itemize}
\end{knowledgebox}

\subsection{Resumen de la sección}

Una nueva contradicción de los productos de Agent es: entre más se parezca la capacidad a «ejecución autónoma de larga duración», más puede empujar al usuario hacia revisiones frecuentes y cambios de contexto. La capa de producto puede aliviar esto mediante el diseño de retroalimentación, pero lo fundamental sigue siendo la capacidad del modelo en tool use, decision making y juicio de tareas. A largo plazo, el usuario no debería entrenarse para adaptarse a un Agent ineficiente; el Agent debería aprender a reducir la carga cognitiva del usuario.

\section{Aprendizaje continuo y memoria: el modelo no debe morir después de salir a producción}

La segunda mitad de la conversación se desplaza hacia preguntas de investigación y la densidad aumenta con claridad. Xie Tianbao (谢天宝) plantea que, durante el último año, buena parte del trabajo en GUI Agent se ha parecido más a tapar huecos de ingeniería: completar datos, completar infra, completar evaluation, completar el entorno. Si se busca una dirección con un cambio de paradigma más profundo, online learning, continual learning, test-time training y self-evolving podrían ser los problemas importantes de uno o dos años.

El contexto de este juicio es que el paradigma actual de entrenamiento de modelos sigue siendo muy «fuera de línea»: preentrenamiento, SFT, RL y después la inferencia en producción. Una vez en producción, el modelo se mantiene básicamente sin cambios. Aunque cada día haya una cantidad enorme de interacciones con usuarios, el modelo no obtiene de ellas, en tiempo real, ninguna ganancia de capacidad. A lo sumo hay retorno de registros, limpieza fuera de línea, reentrenamiento y un nuevo despliegue. El problema es que el ruido de los datos reales de producto es enorme, la retroalimentación de los usuarios es escasa y subjetiva, y convertirla directamente en datos de entrenamiento no es sencillo.

El online RL de Cursor se menciona como una señal: cuando el escenario de producto tiene retroalimentación de alta frecuencia, un objetivo claro y una infrastructure de despliegue rápido, el modelo puede entrar más rápido en el ciclo de «operar el producto es aprender». Pero GUI Agent es mucho más difícil que tab prediction. Se enfrenta a tareas de largo alcance, un espacio de acciones abierto, llamadas a múltiples herramientas, reward escaso y una atribución de fallos compleja.

\begin{importantbox}{El modelo no debe ser solo una herramienta de inferencia después de salir a producción}
Si el modelo interactúa a diario con una gran cantidad de usuarios, pero no puede actualizar a partir de esas interacciones su conocimiento, sus preferencias, su política o su memoria, entonces el proceso de inferencia solo consume cómputo, sin generar un interés compuesto de capacidad. El objetivo del aprendizaje continuo es que la interacción real se convierta en un activo del sistema, y no en basura de registros.
\end{importantbox}

La discusión sobre memory es el núcleo de este segmento. Gu Yu (谷雨) la divide en tres preguntas: cómo se representa la memory, cómo se actualiza y cómo se usa. Esta división es clave, porque muchos productos simplifican memory a «meter el historial en el contexto» o «hacer una base vectorial», pero la memoria de largo plazo real va mucho más allá de la recuperación. Tiene que decidir qué información vale la pena conservar a largo plazo, cómo comprimirla, cómo permitir que el usuario la edite, cómo hacer que las tareas la invoquen y cómo evitar que la información antigua contamine las decisiones nuevas.

Los invitados van y vienen entre la memoria paramétrica y la memoria discreta. La memoria paramétrica tiene un costo de actualización alto, poca plasticidad y es difícil de auditar; la memoria discreta, como texto, SOP, workflow o knowledge graph, es más editable y más fácil de revisar, pero también es propensa a la pérdida de información, a una estructura rígida y a dificultades de context management. Shang Yanyi (尚晏仪) se inclina por representar mucho conocimiento en texto, porque el texto puede entrar en la colaboración organizacional; Xie Tianbao (谢天宝), en cambio, advierte que el cerebro humano en sí mismo no es un sistema puramente simbólico y que, al final, quizá haya que volver a una forma sub-symbolic. Esto no es un desacuerdo, sino una muestra de que el sistema futuro probablemente sea híbrido.

\begin{knowledgebox}{Las tres preguntas de memory}
\begin{tabularx}{\textwidth}{>{\raggedright\arraybackslash}p{0.18\textwidth} Y}
\toprule
Pregunta & Significado clave \\
\midrule
Cómo se representa & Parámetros, texto, vectores, SOP, workflow y knowledge graph tienen distinta plasticidad, distinta auditabilidad y distinto costo de invocación. \\
Cómo se actualiza & Qué se extrae de la interacción, quién lo confirma, cómo se elimina una memoria errónea, cómo un cambio de preferencia sobrescribe un registro antiguo. \\
Cómo se usa & Si es para recuperación, planeación, personalización, control de permisos, recuperación de errores, o para entrar a la siguiente ronda de entrenamiento. \\
\bottomrule
\end{tabularx}
\end{knowledgebox}

La discusión sobre test-time training e intrinsic reward apunta a otra posibilidad: si el modelo puede obtener una señal, durante la inferencia, a partir de su propia distribución, de confidence, de la ruta de muestreo o del resultado de la interacción. La intuición aquí es que el modelo ya absorbió una gran cantidad de conocimiento durante el preentrenamiento y el post-training, y que la respuesta correcta quizá esté en la distribución, solo que el decoding o la búsqueda no la extrajeron. Esta dirección todavía está lejos de madurar, pero vuelve a abrir la frontera entre «inferencia» y «aprendizaje».

\begin{warningbox}{El contexto largo no es memoria a largo plazo}
Meter todo el historial en el context genera problemas de costo, ruido, privacidad, recuperación y consistencia. La memoria de largo plazo debe poder comprimirse, editarse, olvidarse y auditarse; de lo contrario, no es un activo, sino una deuda cognitiva cada vez más pesada.
\end{warningbox}

\subsection{Resumen de la sección}

La discusión sobre el aprendizaje continuo desplaza a GUI Agent de una herramienta de un solo uso hacia un sistema de largo plazo. Si el modelo, después de salir a producción, no puede acumular capacidad a partir de la interacción, no puede formar un interés compuesto real. Memory no es un registro de conversación; incluye, como mínimo, tres cosas: representación, actualización y uso. Es más probable que el futuro sea una mezcla de memoria paramétrica, memoria discreta, workflow y test-time adaptation, y no un único contexto largo o una única base vectorial.
\section{Conocimiento, símbolos y bibliotecas de semántica organizacional}

Este segmento sobre el conocimiento parece filosófico, pero en realidad es muy de ingeniería. Gu Yu (谷雨) menciona que el conocimiento puede dividirse, al menos de manera aproximada, en conocimiento factual, conocimiento de procedimiento y conocimiento situacional. El conocimiento factual responde «qué es», el conocimiento de procedimiento responde «cómo se hace» y el conocimiento situacional registra «qué ocurrió». Para un Agente, lo más difícil casi nunca es lo factual, sino el proceso y la situación: cómo tramitar una visa, cómo hacer un reembolso dentro del sistema corporativo, cómo modificar código siguiendo las convenciones de un equipo determinado, cómo organizar un viaje según las preferencias de un usuario específico.

Esto lleva al tema de workflow y neuro-symbolic. Si se deja que el modelo haga plan en cada paso, el sistema es flexible pero inestable; si el proceso se escribe como código, SOP o RPA, el sistema es estable pero rígido. Un sistema realmente utilizable probablemente necesite convertir uno en otro de manera constante: el modelo abstrae, a partir de la interacción, documentation, SOP, workflow y knowledge graph, y luego, en tiempo de ejecución, esas estructuras restringen al modelo. La estructura simbólica no es un residuo de una época anterior; lo que aporta es consistencia, auditabilidad y coordinación organizacional.

\begin{importantbox}{El nuevo contexto de neuro-symbolic}
Retomar hoy neuro-symbolic no significa volver a las reglas escritas a mano en lugar del modelo, sino hacer que el modelo neuronal abstraiga, en un entorno abierto, estructuras simbólicas editables, verificables y reutilizables. El modelo se encarga de entender la realidad ambigua; la estructura simbólica se encarga de estabilizar el proceso y la semántica compartida.
\end{importantbox}

Después, la discusión sobre la deriva conceptual trae el problema del modelo de vuelta a la organización. Palabras como AGI, Agent, Knowledge y AI-native son usadas una y otra vez por personas y modelos, y van perdiendo una definición precisa. Lo más peligroso es que el modelo reproduce esas palabras de moda hacia las personas, y las personas vuelven a usar esas palabras para tomar decisiones, formando un ciclo de deriva conceptual. Muchos equipos creen estar discutiendo el mismo concepto, pero en realidad cada quien tiene en mente un objeto distinto.

Shang Yanyi (尚晏仪) plantea el problema de manera muy orientada al producto: para una empresa que empieza de cero, lo más fundamental es la biblioteca de semántica. Lo que la empresa considera que es «viajar», qué es «reservar un boleto de avión», qué es «la necesidad del usuario»: estas definiciones son las que, con el tiempo, dan lugar al producto, al proceso y al lenguaje organizacional. Si estas definiciones no están sincronizadas, la coordinación organizacional se vuelve difícil; si las herramientas de IA absorben y generan de manera continua texto interno sin calibrar, a largo plazo pueden incluso deteriorar la racionalidad de la organización.

\begin{warningbox}{Las palabras de moda pueden robar el criterio}
Si palabras como «AGI», «Agent», «AI-native» y «knowledge» no se redefinen, dejan de ser herramientas para pensar y se convierten en sustitutos del pensamiento. El equipo necesita preguntarse primero: ¿a qué nos referimos exactamente en este escenario? ¿cuál es el criterio de éxito? ¿a qué acciones, datos, procesos y límites corresponde?
\end{warningbox}

Xie Tianbao (谢天宝) agrega que la ventaja del mundo symbolic es que sus definiciones son claras y su consistencia es alta. Si se logra conectar symbolic y neural, muchos problemas de consistencia pueden resolverse primero en la capa symbolic, y después comprimirse de vuelta hacia la capa neural o subsymbolic. Esta idea, de hecho, se corresponde con el problema de producto de toda la conversación: el Agente no solo necesita un modelo más fuerte, también necesita límites de acción, definiciones semánticas y memoria organizacional más claros.

\subsection{Resumen de la sección}

El conocimiento no es solo una biblioteca de hechos; incluye también proceso, situación, preferencias y semántica organizacional. Si los sistemas de Agente van a integrarse de manera duradera en organizaciones reales, deben resolver la combinación neuro-symbolic: el modelo se encarga de entender y generalizar, la estructura simbólica se encarga de la consistencia y la coordinación. La deriva conceptual no es un purismo del lenguaje, sino un riesgo de producto y de organización; la biblioteca de semántica se convertirá en un tipo de infraestructura para las empresas de Agentes.
\section{El foso de las startups: qué queda después de la velocidad}

La última parte analiza el foso de las startups. En apariencia se trata de experiencia de emprendimiento, pero en realidad continúa el «problema de la acumulación» planteado antes. Se suele decir que la ventaja de las empresas emergentes es la velocidad, pero Shang Yanyi (尚晏仪) señala que la velocidad no es una lógica completa. El valor de la velocidad está en acumular antes ciertas cosas, y son esas cosas las que pueden llegar a convertirse en un foso.

Ella entiende la competencia temprana de los productos de IA a partir de los puntos de inflexión en la capacidad de los modelos. Cuando un modelo fundacional cruza de pronto cierto umbral crítico, si una empresa pequeña es la más rápida en encontrar un escenario y lanzar un producto, puede aprovechar el vacío que se abre. Ese vacío genera notoriedad, marca y posicionamiento en la mente del usuario. La estructura de costos de difusión en la era de la IA también es distinta de la compra de tráfico de la internet tradicional: si una experiencia realmente rompe la barrera psicológica del usuario, los usuarios la difunden por su cuenta.

Pero la notoriedad no basta. El ejemplo de la industria de viajes muestra que, después de la experiencia de producto, la competencia pasa a la cadena de suministro. Si una aplicación se lanza pronto y tiene buena experiencia, pero el inventario de hoteles es insuficiente, los precios son inestables y el cumplimiento es débil, los usuarios terminan por volver a la plataforma con una cadena de suministro más fuerte. La competencia a largo plazo se convierte en una guerra integral: producto, marca, cadena de suministro, comportamiento del usuario, eficiencia organizacional y capacidad de reinversión forman juntos un círculo virtuoso.

\begin{importantbox}{La velocidad no es el foso; lo es el activo que la velocidad genera}
La velocidad por sí sola no defiende contra quienes llegan después. Solo cuando se transforma en marca, hábito del usuario, cadena de suministro, datos, comunidad, artifact o capacidad organizacional puede llegar a formar una barrera de largo plazo. De lo contrario, en cuanto se difunde la capacidad del modelo fundacional, la ventaja funcional se borra con rapidez.
\end{importantbox}

Xie Tianbao (谢天宝) usa la búsqueda de Google como analogía: si técnicamente toda la internet puede ser indexada, ¿por qué a quienes llegan después les cuesta tanto alcanzar a Google? Posiblemente porque haberlo hecho primero trajo consigo un index más completo, un comportamiento de query más rico, una iteración de producto más rápida y un hábito de marca más fuerte. Esta analogía nos recuerda que la ventaja de moverse primero no es algo místico, sino la acumulación de un conjunto concreto de artifacts.

También discutieron qué tipo de acumulación es más difícil de que las grandes empresas absorban. Los datos por sí solos pueden ser igualados con más cómputo y más usuarios; pero los artifacts de escenario, los entornos AI-native, las relaciones de cadena de suministro, la confianza de la comunidad, la semántica organizacional y los detalles del producto son más difíciles de replicar. Gu Yu (谷雨) menciona que NeoCognition quiere mejorar el entorno digital mismo, para que el entorno se vuelva más AI-native; si el resultado es justamente el artifact del entorno ya mejorado, entonces la acumulación de haberlo hecho antes es el producto mismo, no datos de entrenamiento indirectos.

\begin{knowledgebox}{Varios tipos de activos acumulables de las startups de IA}
\begin{tabularx}{\textwidth}{>{\raggedright\arraybackslash}p{0.23\textwidth} Y}
\toprule
Tipo de activo & Por qué puede formar una barrera \\
\midrule
Marca y posicionamiento mental & Cuando surge una nueva necesidad, el usuario piensa primero en ti, lo que reduce el costo posterior de adquisición de usuarios. \\
Cadena de suministro y cumplimiento & En los escenarios transaccionales, la experiencia termina por hacerse realidad a través del inventario, el precio, el servicio y los límites de responsabilidad. \\
Artifact de escenario & Si el resultado mismo es un entorno, workflow o estructura de conocimiento reutilizable, haberlo hecho temprano deja sedimentado un activo directo. \\
Comunidad y confianza & Las empresas pequeñas pueden generar retención mediante relaciones más cercanas con los usuarios; esta conexión no se sustituye fácilmente por la mera funcionalidad. \\
Repositorio de semántica organizacional & La calibración de largo plazo de los conceptos del escenario, los procesos y las necesidades del usuario eleva la eficiencia de la iteración. \\
\bottomrule
\end{tabularx}
\end{knowledgebox}

Al final, Shang Yanyi enfatiza la conexión entre las personas. Si un producto es solo una herramienta, el hábito del usuario puede migrar con rapidez; pero si el producto implica que el usuario entra a una comunidad, confirma cierta identidad y establece relaciones con un grupo de personas, la retención será mayor. Esta es una oportunidad importante de las empresas pequeñas frente a las grandes: no ganar en todos los recursos, sino construir una red de confianza más densa dentro de un grupo concreto de personas.

\subsection{Resumen de la sección}

El foso no es un sustantivo estático, sino una estructura de interés compuesto dentro de una competencia dinámica. La velocidad temprana de las startups de IA debe transformarse en activos acumulables; el resultado de la etapa media y tardía depende de si el producto, la cadena de suministro, la comunidad, los artifacts y la eficiencia organizacional logran formar un círculo virtuoso. Cuanto más fácil sea que una plataforma de modelo fundacional absorba una capacidad, menos apta será como barrera de largo plazo; cuanto más cercano esté un activo al escenario y a la relación con el usuario, más probable será que perdure.
\section{Síntesis: la siguiente parada del GUI Agent no es dar clics con más destreza}

Toda la conversación termina por resumirse en un juicio: la siguiente parada del GUI Agent no es dar clics con más destreza en la pantalla, sino insertarse mejor en la vida digital real. Una vez que entra al mundo real, los problemas se vuelven de inmediato más densos: cómo se evalúa, cómo se cobra, cómo se reduce la carga del usuario, cómo se usa el contexto, cómo se aprende de manera continua, cómo se mantiene la semántica de la organización, cómo se forma un activo acumulable. Ninguno de estos problemas se resuelve con una sola demostración.

Si se comprimen los puntos de vista de los tres invitados en una estructura, se obtiene la siguiente tabla:

\begin{longtable}{>{\raggedright\arraybackslash}p{0.24\textwidth} >{\raggedright\arraybackslash}p{0.68\textwidth}}
\toprule
Enfoque del problema & Juicio central en la conversación \\
\midrule
\endfirsthead
\toprule
Enfoque del problema & Juicio central en la conversación \\
\midrule
\endhead
La esencia del GUI Agent & No es una herramienta que «mira la pantalla y hace clic», sino un sistema agente que convierte el objetivo del usuario en acciones dentro del entorno digital. \\
La revelación de Operator & Demuestra que la dirección es válida, pero también expone que un «proceso de operación visible» no equivale a una buena experiencia de producto. \\
Los límites del benchmark & El benchmark es una señal, no una respuesta final; el producto real necesita calibrarse con retroalimentación del usuario y resultados comerciales. \\
La dificultad de la recompensa & El resultado final suele estar demasiado lejos y la retroalimentación subjetiva es demasiado ruidosa; el producto de Agent necesita experimentación continua y atribución de fallas. \\
Modelo de negocio & El Agent se parece más a un trabajador intelectual y a un agente transaccional; la suscripción y la comisión pueden resultar más naturales que la publicidad. \\
UDA y el espacio de acciones & Deep Research, Coding, Browser y GUI no deberían fragmentarse; un Agent maduro debe elegir la acción en función del objetivo. \\
Experiencia en tareas largas & Poder ejecutarse durante mucho tiempo no es el valor en sí; reducir la espera, la verificación y el cambio de contexto es la dirección real de mejora de la experiencia. \\
UI generativa & Lo valioso es un componente estable con orquestación dinámica, no reinventar la interfaz cada vez. \\
Aprendizaje continuo & El modelo, una vez en producción, no debería ser solo una herramienta de inferencia; la interacción real debería volverse fuente de capacidad y de memoria. \\
Memory & Lo clave son tres cosas: representación, actualización y uso; no meter el historial dentro del context. \\
Deriva conceptual & Las palabras de moda contaminan el juicio de la organización; la empresa necesita una biblioteca semántica para mantener definiciones comunes. \\
Foso defensivo de la startup & La velocidad no es un foso defensivo; lo que sí lo es son la marca, la cadena de suministro, la comunidad, los artifact y la capacidad organizacional que la velocidad genera. \\
\bottomrule
\end{longtable}

Lo mejor de este episodio es que no se queda en el optimismo simplista de «el Agent va a cambiar todo muy pronto». Al contrario, muestra una y otra vez una realidad: cuando el Agent de verdad va a entrar en la vida del usuario y en el proceso de la empresa, la capacidad técnica es solo la primera capa. Después vienen el mecanismo de evaluación, la distribución comercial, la carga de interacción, el límite de privacidad, el aprendizaje continuo, la coherencia semántica y el interés compuesto organizacional. El futuro del GUI Agent tal vez no pertenezca al sistema que mejor demuestre clics de mouse, sino al que logre organizar, paso a paso, estos problemas densos.
\section*{Apéndice: índice de juicios clave}

\begin{longtable}{>{\raggedright\arraybackslash}p{0.34\textwidth} >{\raggedright\arraybackslash}p{0.58\textwidth}}
\toprule
Juicio clave & Cuestión que aborda \\
\midrule
\endfirsthead
\toprule
Juicio clave & Cuestión que aborda \\
\midrule
\endhead
\textit{Deep Research, Coding y GUI Agent no pueden hacerse por completo por separado, porque todos ocurren en la misma computadora y en el mismo espacio de objetivos.} & Este es el punto de partida intuitivo del Unified Action Space. El objetivo del usuario cruza herramientas de manera natural, y las fronteras del producto no deberían obligar al usuario a trasladar el contexto. \\
\textit{Operator hizo que se reconociera la dirección, pero no necesariamente es un buen producto.} & Esto distingue la prueba técnica de la experiencia del usuario. Que se pueda mostrar a la IA operando no significa que el usuario quiera ver a la IA operar lentamente. \\
\textit{Las puntuaciones de los benchmarks se parecen cada vez más a una señal que a una meta final.} & Esto indica que la evaluación está pasando de una comparación puramente académica a una calibración con escenarios reales. Las puntuaciones siguen siendo útiles, pero los límites de su interpretación deben ser más claros. \\
\textit{El reward final de un producto de viajes es vender un boleto, pero ese reward está demasiado lejos del entrenamiento del modelo.} & Esto revela el problema del reward en los agentes industriales: el resultado final es claro, pero la atribución del proceso es ambigua y la retroalimentación del usuario es subjetiva y muy ruidosa. \\
\textit{Lo que produce el agente es trabajo intelectual, no solo distribución de información.} & Este es el núcleo de la discusión sobre el modelo de negocio. La publicidad vende atención; el agente se parece más a un secretario, un consultor, un intermediario o un servicio de representación. \\
\textit{Poder realizar una tarea de una hora no es bueno en sí mismo si la misma calidad puede lograrse en 30 segundos; nadie quiere esperar una hora.} & Esto critica el mercadeo de las tareas largas. La capacidad de ejecución autónoma debe traducirse en una carga menor para el usuario, no en una espera más larga. \\
\textit{La memoria debe abordar al menos tres cuestiones: representación, actualización y uso.} & Esto eleva la memoria de largo plazo de «historial de conversación» a un problema de diseño de sistemas. La memoria debe poder editarse, olvidarse y reutilizarse. \\
\textit{Las palabras de moda se propagan una y otra vez entre modelos y personas, y al final producen una deriva de concepto.} & Esto convierte el problema del lenguaje en la era de la IA en un problema organizacional. Sin una biblioteca semántica, la empresa pierde el juicio compartido en medio de las palabras de moda. \\
\textit{La velocidad no es un foso; lo que importa es lo que se acumula después de la velocidad.} & Este es el núcleo del apartado sobre startups. La velocidad solo tiene valor de largo plazo cuando se convierte en marca, cadena de suministro, comunidad, artifact o capacidad organizacional. \\
\bottomrule
\end{longtable}

\end{document}
